\chapter{Introdução}
\label{cap:intro}

A Lei de Acesso à Informação (LAI), promulgada em 2011, consolidou o direito do cidadão de obter dados produzidos pelo Estado brasileiro \cite{Lei:12.527:2011}. Como desdobramento dessa legislação, os portais municipais de transparência passaram a publicar informações sobre orçamento, finanças e administração. Publicar dados, por si só, não basta. A linguagem técnica do orçamento, a fragmentação das informações em múltiplas tabelas e a ausência de ferramentas de análise mantêm o cidadão distante da compreensão de como os recursos públicos são gastos \cite{matiaspereira2010,melo2019proposta}.

O Data Warehouse (DW) surge como resposta tecnológica a esse problema. Trata-se de um repositório que integra dados de fontes heterogêneas e os organiza para análise \cite{kimball2013data}. A modelagem dimensional, consolidada por \citeonline{kimball2013data}, prioriza a usabilidade da consulta em relação à otimização das transações. Em vez de planilhas que exigem do usuário cruzar referências manualmente, tabelas de fatos e dimensões permitem que perguntas relevantes sejam respondidas com poucos cliques: "Quanto Juiz de Fora investiu em educação nos últimos três anos?", "Qual o tempo médio de pagamento aos fornecedores da área de saúde?", "Há concentração de recursos em poucos favorecidos?". Essas perguntas, antes restritas a especialistas em finanças públicas, passam a ser acessíveis a jornalistas, conselheiros municipais e organizações da sociedade civil \cite{marco2022transparencia}.

A transparência orçamentária brasileira convive com um paradoxo. A base legal existe, os portais estão no ar, os dados são publicados. Mesmo assim, a participação cidadã na fiscalização do gasto público permanece restrita. A literatura aponta três lacunas que sustentam essa distância entre a disponibilidade formal e o uso real da informação: a lacuna social, a técnica e a pedagógica.

A lacuna social refere-se às barreiras ao letramento financeiro e digital. Cidadãos comuns têm dificuldade não só em acessar dados, mas também em relacioná-los à própria experiência. O orçamento público, apresentado em categorias contábeis como "elemento de despesa", "função programática" e "fonte de recurso", permanece distante da realidade da escola, do posto de saúde e da via pública. A opacidade técnica desmobiliza a participação e reforça o privilégio de quem tem formação especializada ou acesso a intermediários, como ONGs e jornalismo investigativo. O resultado é um déficit democrático: a transparência formal coexiste com a opacidade prática.

A lacuna técnica está nos próprios portais. Os dados costumam ser disponibilizados em formatos pouco amigáveis à análise, como PDFs sem estrutura ou planilhas sem padronização, e raramente há APIs ou recursos para o cruzamento de informações entre bases distintas (despesas, receitas, licitações). Quem deseja produzir análises temporais ou identificar padrões anômalos enfrenta um trabalho manual de coleta, limpeza e consolidação que poucos podem custear em termos de tempo e de especialização. O Data Warehouse altera essa equação ao consolidar o esforço analítico em uma estrutura reutilizável.

A lacuna pedagógica é a menos debatida e fundamenta-se na obra de Paulo Freire, em especial na noção de educação libertadora e no conceito de conscientização \cite{freire1967liberdade}. Freire argumenta que a transformação social depende da capacidade de "ler o mundo" criticamente, superando a passividade diante das estruturas de dominação. Aplicado ao orçamento público, isso significa que a democratização do poder vai além da publicação passiva de dados: exige instrumentos que capacitem o cidadão a formular perguntas, testar hipóteses e construir suas próprias leituras sobre a gestão pública \cite{freire1996autonomia}. O conceito freireano de "inédito viável" descreve com precisão o lugar do Data Warehouse civil. Trata-se de uma forma de participação ainda incipiente, em que o cidadão questiona não só quanto foi gasto, mas também por que aquela área recebeu mais recursos, se o investimento responde às necessidades locais e se há sinais de favorecimento a fornecedores específicos.

A escolha da execução orçamentária municipal como objeto de estudo não é arbitrária. É na esfera local que as políticas públicas ganham forma concreta no cotidiano dos moradores, e é onde o controle social direto encontra mais espaço \cite{novaes2014orcamento}. Ao documentar a construção de um DW para a gestão municipal, o trabalho oferece uma referência metodológica que outros municípios poderiam adaptar, sujeita aos ajustes específicos de parsing de cada portal. O recorte também preenche uma lacuna específica na interdisciplinaridade. Estudos sobre transparência pública costumam permanecer no terreno da análise política ou jurídica. Estudos sobre DW raramente discutem implicações cívicas. Esse trabalho se posiciona no encontro entre as duas e sustenta uma tese: a arquitetura de dados é uma escolha política, pois define quem pode participar do debate público e em que termos.

Metodologicamente, o trabalho adota \textit{Design Science Research} (DSR) como abordagem orientadora \cite{hevner2004, dresch2015}. A escolha responde à natureza do problema: o produto central da pesquisa é um artefato (o Data Warehouse e o painel analítico associado) cuja relevância só pode ser aferida em relação ao problema que ele endereça. DSR articula, em um único ciclo metodológico, a construção e a avaliação desse artefato, evitando o descolamento entre a engenharia da solução e a verificação de sua utilidade. O Capítulo~\ref{cap:metodologia} detalha as fases do ciclo adotado e a correspondência entre cada uma e os capítulos do trabalho.

\section{Problema de Pesquisa}

A pergunta que orienta este trabalho pode ser formulada de modo direto: \textit{Como organizar, em um modelo dimensional, os dados publicados pelo Portal da Transparência da Prefeitura de Juiz de Fora, de forma a reduzir as barreiras técnicas que hoje afastam cidadãos não especialistas da execução orçamentária municipal?} A resposta tem dois lados. 

Do lado técnico, a pergunta se desdobra em escolhas concretas de modelagem: qual a granularidade adequada para os fatos, quais dimensões refletem as classificações oficiais sem reproduzir a barreira semântica que elas impõem, como tratar inconsistências da fonte sem invisibilizá-las e como garantir reprodutibilidade do \textit{pipeline}. Do lado conceitual, a pergunta exige justificar por que essa reorganização importa, ou seja, em nome de quem e para que cidadão ela é feita. Os capítulos seguintes desenvolvem as duas frentes em paralelo, sustentando a ideia de que a separação entre técnica e propósito empobrece os projetos de transparência.

\section{Objetivos}
\subsection*{Objetivo Geral}
Desenvolver um DW para análise da execução orçamentária do município de Juiz de Fora, reorganizando dados públicos dispersos em uma estrutura dimensional acessível que facilite a interpretação e a fiscalização cidadã dos gastos públicos.

\subsection*{Objetivos Específicos}

\begin{itemize}
    \item Coletar e documentar dados de receitas e despesas do Portal da Transparência da Prefeitura de Juiz de Fora, identificando inconsistências nas fontes disponíveis.
    \item Modelar dimensionalmente a execução orçamentária segundo a metodologia Kimball, mapeando as classificações oficiais (MCASP/MTO) em fatos e dimensões.
    \item Implementar processos de ETL para extração, limpeza, transformação e carga de dados, tratando heterogeneidades de formato e de nomenclatura.
    \item Construir dashboards interativos que traduzam classificações orçamentárias oficiais em visualizações acessíveis, com indicadores derivados das tabelas de fatos do warehouse.
    \item Avaliar, em tarefas representativas de fiscalização, o painel desenvolvido em comparação com o Portal da Transparência da Prefeitura de Juiz de Fora, identificando ganhos e limitações de usabilidade analítica.
\end{itemize}

\section{Delimitação do Escopo}
\label{sec:escopo}

Este trabalho concentra-se na execução orçamentária do município de Juiz de Fora, com escopo deliberadamente restrito a quatro dimensões.

A primeira dimensão é geográfica. A aplicação do modelo dimensional proposto restringe-se ao Portal da Transparência da Prefeitura de Juiz de Fora. A replicação em outros municípios é discutida como potencial metodológico, com base na arquitetura em camadas adotada, mas não é validada empiricamente no escopo deste trabalho.

A segunda dimensão é temática. O modelo cobre receitas e despesas no nível em que são publicadas pelo portal municipal, complementadas pelo Ementário de Classificação por Natureza da Receita da Secretaria do Tesouro Nacional e pela classificação funcional da Lei Orçamentária Anual do exercício corrente. Licitações, contratos, folha de pessoal e dados patrimoniais não são integrados ao warehouse, ainda que reconhecidos como extensões naturais do modelo e registrados como trabalhos futuros.

A terceira dimensão refere-se à granularidade. A fonte publica relatórios mensais consolidados e não transações individuais. O modelo dimensional respeita o grão da fonte, sem simular atomicidade inexistente. Análises de empenho a empenho ou item a item exigiriam bases distintas, fora do escopo da publicação atual do portal.

A quarta dimensão é de avaliação. A verificação do artefato é realizada por meio de uma análise funcional comparativa entre o painel desenvolvido e o portal original, em tarefas representativas de fiscalização. Testes de usabilidade com usuários finais, embora reconhecidos como um caminho natural de continuidade, não integram o escopo deste trabalho, por exigirem um protocolo de pesquisa com seres humanos e um prazo de execução incompatíveis com o cronograma de um Trabalho de Conclusão de Curso.

\section{Estrutura do Trabalho}

O restante deste texto está organizado em sete capítulos. O Capítulo~\ref{cap:fundamentacao} apresenta os conceitos teóricos que sustentam o trabalho: orçamento público brasileiro, controle social, Data Warehouse e modelagem dimensional. O Capítulo~\ref{cap:trabalhos} discute trabalhos relacionados às três interseções relevantes (DW aplicado ao setor público, avaliação de portais de transparência e visualização de dados orçamentários) e posiciona esse trabalho em relação ao que já existe. O Capítulo~\ref{cap:metodologia} descreve a metodologia adotada. O Capítulo~\ref{cap:dw} documenta a arquitetura, o modelo dimensional e o pipeline de ETL implementados. O Capítulo~\ref{cap:resultados} apresenta os \textit{dashboards} e a análise dos dados de Juiz de Fora. O Capítulo~\ref{cap:avaliacao} avalia a usabilidade do sistema e discute suas limitações e implicações democráticas. O Capítulo~\ref{cap:conclusao} sintetiza as contribuições e aponta direções para trabalhos futuros.
